\newpage
\setcounter{table}{0}
\setcounter{figure}{0}
\section{总结与展望}

本章的内容主要包含两部分。具体地，本章在~\ref{summarize}节中对本文的研究工作进行了总结，在~\ref{future}节中对本文工作的局限性进行了分析和探讨，为未来的研究工作奠定基础。


\subsection{总结}\label{summarize}





总的来说，本文的主要研究内容如下所示：
\begin{itemize}
	\item[(1)]  \textbf{设计并实现了异构审计日志解析与流式事务重组机制}
	
    针对传统基于网络数据包捕获的回放方案中难以获取准确事务边界以及对多客户端会话解析混乱的问题，本文提出了一种基于插件化架构的异构审计日志解析机制。为了高效处理海量且乱序的日志流，设计了基于哈希映射的单遍流式事务重组算法，在 $O(N)$ 时间复杂度内实现了会话事件的精确分流和时序上下文的无缝整合，为后续的高保真回放奠定了坚实的数据基础。
	
	\item[(2)] \textbf{构建了高并发、时序对齐的流量回放引擎}
	
    为了在回放环境中忠实重现生产负载的高并发特性和时序特征，本文基于 Go 语言的 Goroutine 机制构建了“每会话一协程”的回放调度架构。此外，创新性地提出了一种基于滑动时间窗口的时序对齐机制，通过精确维持事务之间的执行间隔，不仅保障了回放的真实性，同时赋予了系统支持任意倍速调整的灵活性，有效降低了因系统底层调度引发的时序偏离。
	
	\item[(3)] \textbf{设计了多维度的回放保真度评估体系与实验验证}

    通过建立包括功能一致性、吞吐量保真度、时延分布保真度和调度稳定性在内的多维度性能评估体系，克服了以往系统难以量化回放质量的短板。基于 TPC-C 50GB 规模的综合实验证明，系统在支撑极限达 4.5 万 QPS 的负载下，不仅保持了 99.99\% 的 SQL 执行成功率，而且在模拟复杂事务的性能趋势和响应尾延分布时，与源端生产基准的保真度均超过 98\%，体现出极高的吞吐与调度稳定性。
	
\end{itemize}


\subsection{展望}\label{future}
尽管本文提出并实现了一种高保真的基于审计日志的流量回放系统，但在实际复杂的生产环境部署和系统通用性方面仍有进一步完善的空间。未来的研究与优化方向主要集中在以下几个方面：

\begin{itemize}
	\item[(1)]  \textbf{支持更多异构数据库的审计日志格式}
	
    当前的解析重构体系主要基于 PostgreSQL 的 pgaudit 插件生成的格式进行深度优化与测试验证。考虑到现代企业 IT 架构的多样化，未来可通过扩展解析器插件（Parser Plugin），使系统支持如 MySQL 的原生审计插件、Oracle 的 Database Audit 相关格式，从而进一步提升流量回放系统的普适性和跨数据库平台的迁移测试能力。
	
	\item[(2)] \textbf{深化回放差异的自动化归因与智能诊断}
	
    现有系统虽然具备精确捕捉并报告源端与目标端在结果集条数及错误状态码上的差异，但针对复杂场景下引起这些差异的根本原因仍需人工介入分析。未来可引入机器学习或规则引擎模型，对触发差异的慢查询执行计划、锁竞争事件进行自动收集和归类，提升根因定位的智能化水平，大幅降低变更测试的排查成本。
	
	\item[(3)] \textbf{探索基于历史日志特征的智能负载突变与生成}
	
    本文的回放引擎严格忠实于捕获的历史时序与真实日志。然而，在进行极端的容量规划压测或防御性演练时，系统常常需要面对超越历史记录中的异常突增或新型访问模式。未来研究可考虑结合马尔可夫链或其他生成模型，从历史审计日志中学习 SQL 模版特征和参数分布，衍生合成出“形似且加强”的变异负载，进而针对目标数据库开展更加前瞻性和探索性的系统压力测试。
	
\end{itemize}

